\section{End-to-End Integrations}
\label{sec:integrations}

This section focuses only on the dense validation pair:
a systematic dense banded outer together with the dense time-varying recursive inner.
The proof is a direct first-moment argument.
The outer contributes an expected spectrum envelope, while the inner contributes two analytic tails:
one for low/intermediate weights and one geometric tail above a logarithmic cutoff.

\subsection{Direct first-moment reduction}
\label{sec:integration-dense-dense}

Take as outer code the systematic dense log-memory construction of
Section~\ref{sec:outer-dense-logmem}, with memory
\[
M=\lceil c\log_2 n\rceil,
\]
and as inner code the dense time-varying recursive construction of
Section~\ref{sec:inner-dense}, with memory
\[
m=\lceil \gamma \log_2 n\rceil.
\]
The serial concatenation still has one random interleaver and \(O(\log n)\) taps per symbol in both stages,
hence \(O(n\log n)\) encoding work.

\begin{theorem}[Direct dense+dense first-moment bound]
\label{thm:dense-dense-direct-first-moment}
Fix \(\delta\in(0,1/2)\).
Suppose the outer ensemble satisfies
\[
\EE[\enum_{\Cone,h}] \le A_h
\qquad (1\le h\le n),
\]
and the inner satisfies
\[
p_h(\delta)
:=
\Pr[\wt(\Enc_{\mathrm{in}}(U))\le \delta n\mid \wt(U)=h]
\le B_h
\qquad (1\le h\le n),
\]
where \(U\) is uniform on the weight-\(h\) slice of \(\{0,1\}^n\).
Then
\begin{equation}
\label{eq:dense-dense-direct-first-moment}
\Pr[d_{\min}(\Csc)\le \delta n]
\;\le\;
\sum_{h=1}^{n} A_h\,B_h.
\end{equation}
\end{theorem}

\begin{proof}
By the standard first-moment reduction,
\[
\Pr[d_{\min}(\Csc)\le \delta n]
\le
\sum_{h=1}^{n}\EE[\enum_{\Cone,h}]\,p_h(\delta),
\]
because the interleaver makes every fixed weight-\(h\) outer codeword uniform on the weight-\(h\) slice before entering the inner code.
Apply the envelopes \(A_h\) and \(B_h\).
\end{proof}

\subsection{A baseline linear-distance theorem}
\label{sec:integration-dense-dense-baseline}

We now combine the global systematic-outer spectrum envelope with the two dense-inner bounds.

\begin{theorem}[Systematic dense outer + dense inner have linear distance]
\label{thm:dense-dense-integration}
Fix any \(z\in(0,\sqrt2-1)\), and define
\[
a(z):=\frac{1+z}{2}.
\]
Choose the outer memory constant \(c\) so that
\begin{equation}
\label{eq:dense-c-condition}
c>\frac{1}{\log_2(1/a(z))}.
\end{equation}
Let \(\beta_{\mathrm{out}}>1\) be the corresponding outer spectrum constant from
Corollary~\ref{lem:outer-dense-spectrum}.

Next choose constants \(\kappa>0\), \(\rho\in(0,1)\), \(\delta\in(0,1/2)\), and \(\gamma>1\) such that
\begin{equation}
\label{eq:dense-dense-parameter-conditions}
\kappa\log_2(\beta_{\mathrm{out}})<1,
\qquad
2\delta<\rho<\beta_{\mathrm{out}}^{-1},
\qquad
\gamma>1+\kappa\log_2\!\Bigl(\frac{1}{\rho}\Bigr).
\end{equation}
Then the resulting dense outer + dense inner concatenation satisfies
\[
\Pr[d_{\min}(\Csc)\le \delta n]=o(1).
\]
In particular, \(d_{\min}(\Csc)=\Omega(n)\) with high probability.
\end{theorem}

\begin{proof}
Set
\[
h_\star:=\lfloor \kappa\log_2 n\rfloor.
\]
We split the first-moment sum from Theorem~\ref{thm:dense-dense-direct-first-moment} into
\[
S_{\mathrm{low}}
:=
\sum_{h=1}^{h_\star} \EE[\enum_{\Cone,h}]\,p_h(\delta),
\qquad
S_{\mathrm{high}}
:=
\sum_{h=h_\star+1}^{n} \EE[\enum_{\Cone,h}]\,p_h(\delta).
\]

\paragraph{Outer side.}
Corollary~\ref{lem:outer-dense-spectrum} gives
\[
\EE[\enum_{\Cone,h}] \le \beta_{\mathrm{out}}^h
\qquad (1\le h\le n).
\]

\paragraph{Low-weight inner bound.}
Fix any \(\xi>0\) and let
\[
L:=\lceil(2+\xi)\delta n\rceil,
\qquad
b_{\delta,\xi}(n):=
\PP\!\left[\Binomial(L,\tfrac12)\le \delta n\right].
\]
Apply the small-weight envelope \eqref{eq:dense-inner-smallw-envelope} from
Corollary~\ref{cor:dense-inner-piecewise}:
\[
p_h(\delta)
\le
\frac{h(h-1)}{n-1} + (L2^{-m})^h + h\,b_{\delta,\xi}(n)
\qquad (1\le h\le h_\star).
\]
Therefore
\[
S_{\mathrm{low}}
\le
\sum_{h=1}^{h_\star}\beta_{\mathrm{out}}^h\frac{h(h-1)}{n-1}
\;+\;
\sum_{h=1}^{h_\star}\beta_{\mathrm{out}}^h(L2^{-m})^h
\;+\;
b_{\delta,\xi}(n)\sum_{h=1}^{h_\star} h\beta_{\mathrm{out}}^h.
\]

For the first term,
\[
\sum_{h=1}^{h_\star}\beta_{\mathrm{out}}^h\frac{h(h-1)}{n-1}
\le
\frac{h_\star^2}{n}\sum_{h=1}^{h_\star}\beta_{\mathrm{out}}^h
\le
n^{\kappa\log_2(\beta_{\mathrm{out}})-1}(\log n)^{O(1)}
=o(1)
\]
by the first condition in \eqref{eq:dense-dense-parameter-conditions}.

For the second term, since \(m=\gamma\log_2 n\),
\[
L2^{-m}=O(n^{1-\gamma}),
\]
so
\[
\beta_{\mathrm{out}}L2^{-m}=O(n^{1-\gamma})=o(1)
\]
because \(\gamma>1\). Hence
\[
\sum_{h=1}^{h_\star}\beta_{\mathrm{out}}^h(L2^{-m})^h=o(1).
\]

For the third term, \(b_{\delta,\xi}(n)\) is exponentially small in \(n\), while
\[
\sum_{h=1}^{h_\star} h\beta_{\mathrm{out}}^h
\le
(\log n)^{O(1)}\,n^{\kappa\log_2(\beta_{\mathrm{out}})},
\]
so this contribution is also \(o(1)\).
Thus
\[
S_{\mathrm{low}}=o(1).
\]

\paragraph{High-weight inner bound.}
By Corollary~\ref{cor:dense-inner-constant-geometric}, the conditions
\[
2\delta<\rho,
\qquad
\gamma>1+\kappa\log_2(1/\rho)
\]
imply that for all \(h\ge h_\star+1\),
\[
p_h(\delta)\le 3\rho^h.
\]
Therefore
\[
S_{\mathrm{high}}
\le
3\sum_{h=h_\star+1}^{n}(\beta_{\mathrm{out}}\rho)^h.
\]
Since \(\beta_{\mathrm{out}}\rho<1\), this is a geometric tail, and
\[
S_{\mathrm{high}}
\le
O(1)\,(\beta_{\mathrm{out}}\rho)^{h_\star}
=
n^{-\kappa\log_2(1/(\beta_{\mathrm{out}}\rho))+o(1)}
=o(1).
\]

Combining the low- and high-weight windows proves
\[
\Pr[d_{\min}(\Csc)\le \delta n]
\le
S_{\mathrm{low}}+S_{\mathrm{high}}
=
o(1).
\]
\end{proof}

\paragraph{What this baseline theorem buys us.}
Theorem~\ref{thm:dense-dense-integration} is fully analytic and already matches the qualitative picture from the
enumerator experiments:
\begin{itemize}
\item the systematic outer removes collapse to tiny intermediate weights,
\item the outer contributes a global geometric spectrum envelope,
\item the inner handles small weights by the run-aware bound and large weights by a geometric cutoff.
\end{itemize}
This is the baseline dense+dense theorem. The next quantitative question is how close one can push \(\delta\) toward
the random-linear benchmark while remaining within the admissible parameter region
\eqref{eq:dense-dense-parameter-conditions}.

\begin{corollary}[Explicit admissible dense+dense parameter region]
\label{cor:dense-dense-explicit-region}
Fix any \(z\in(0,\sqrt2-1)\), and choose
\[
c>\frac{1}{\log_2(2/(1+z))}.
\]
Let \(\kappa>0\), \(\rho\in(0,z)\), \(\delta\in(0,\rho/2)\), and \(\gamma>1\) satisfy
\[
\kappa\log_2(1/z)<1,
\qquad
\gamma>1+\kappa\log_2\!\Bigl(\frac{1}{\rho}\Bigr).
\]
Then the systematic dense outer with memory \(M=\lceil c\log_2 n\rceil\) concatenated with the dense inner of memory
\(m=\lceil \gamma\log_2 n\rceil\) satisfies
\[
\Pr[d_{\min}(\Csc)\le \delta n]=o(1).
\]
\end{corollary}

\begin{proof}
For all sufficiently large \(n\), Corollary~\ref{lem:outer-dense-spectrum} gives
\[
\EE[\enum_{\Cone,h}] \le z^{-h},
\]
so Theorem~\ref{thm:dense-dense-integration} applies with \(\beta_{\mathrm{out}}=z^{-1}\).
The stated conditions are exactly the theorem's hypotheses rewritten in terms of \(z\).
\end{proof}

\begin{corollary}[One concrete baseline choice]
\label{cor:dense-dense-concrete-choice}
Choose
\[
z=\frac13,
\qquad
\kappa=\frac12,
\qquad
\rho=\frac14.
\]
If
\[
c>\frac{1}{\log_2(3/2)},
\qquad
\gamma>2,
\qquad
\delta<\frac18,
\]
then the systematic dense outer with memory \(M=\lceil c\log_2 n\rceil\) and the dense inner with memory
\(m=\lceil \gamma\log_2 n\rceil\) satisfy
\[
\Pr[d_{\min}(\Csc)\le \delta n]=o(1).
\]
\end{corollary}

\begin{proof}
For \(z=\tfrac13\), we have \(a(z)=\tfrac23\), so the outer condition becomes
\[
c>\frac{1}{\log_2(3/2)}.
\]
Also,
\[
\kappa\log_2(1/z)=\frac12\log_2 3<1,
\]
and
\[
1+\kappa\log_2(1/\rho)
=
1+\frac12\log_2 4
=2.
\]
Finally, \(\delta<1/8\) implies \(2\delta<1/4=\rho<z=1/3\).
Thus Corollary~\ref{cor:dense-dense-explicit-region} applies.
\end{proof}

\begin{remark}[Memory-vs.-blocklength interpretation]
\label{rem:dense-dense-memory-scaling}
The concrete choice in Corollary~\ref{cor:dense-dense-concrete-choice} can be restated directly in terms of the
log-memory parameter.
If the outer memory is
\[
M=\lceil c\log_2 n\rceil
\]
with
\[
c>\frac{1}{\log_2(3/2)}\approx 1.71,
\]
and the inner memory is
\[
m=\lceil \gamma\log_2 n\rceil
\qquad\text{with}\qquad \gamma>2,
\]
then the dense serial concatenation has linear minimum distance with high probability for every \(\delta<1/8\).
Equivalently, fixing a memory budget \(M\), the baseline proof certifies linear distance throughout the regime
\[
n \le 2^{M/c}
\]
for any fixed \(c>1/\log_2(3/2)\).
This does not yet match the stronger finite-length trend suggested by the enumerator, but it gives a rigorous
logarithmic-memory scaling window from which sharper dense bounds can now be pursued.
\end{remark}

\subsection{A linear-weight exponent criterion}
\label{sec:integration-dense-dense-linear}

The next sharpening step is to replace the coarse high-weight geometric tail by a direct comparison of the outer and
inner linear-weight exponents.

\begin{theorem}[Dense+dense linear-weight criterion]
\label{thm:dense-dense-linear-weight-criterion}
Fix constants \(\eta_0\in(0,1)\), \(\delta\in(0,1/2)\), \(\xi>0\), and \(\theta\in(0,1-\eta_0)\).
For \(\eta\in[\eta_0,1-\theta]\), let
\[
\Phi_{\mathrm{out}}^{\mathrm{sys}}(\eta)
\]
be the outer exponent from Theorem~\ref{thm:outer-dense-linear-exponent}, and let
\[
E_{\mathrm{in}}(\eta;\theta,\delta,\xi)
\]
be the dense-inner exponent from Theorem~\ref{prop:dense-inner-linear-exponent}.
Assume there exists \(\tau>0\) such that
\begin{equation}
\label{eq:dense-dense-linear-gap}
\Phi_{\mathrm{out}}^{\mathrm{sys}}(\eta)
\;-\;
E_{\mathrm{in}}(\eta;\theta,\delta,\xi)
\;\le\;
-\tau
\qquad\text{for every }\eta\in[\eta_0,1-\theta].
\end{equation}
Then the total first-moment contribution from outer weights \(h\ge \eta_0 n\) satisfies
\[
\sum_{h=\lceil \eta_0 n\rceil}^{\lfloor (1-\theta)n\rfloor}
\EE[\enum_{\Cone,h}]\,p_h(\delta)
\;\le\;
n^{O(1)}\,2^{-\tau n}.
\]
\end{theorem}

\begin{proof}
For \(h=\lfloor \eta n\rfloor\) with \(\eta\in[\eta_0,1-\theta]\), Theorem~\ref{thm:outer-dense-linear-exponent} gives
\[
\EE[\enum_{\Cone,h}]
\le
n^{O(1)}\,2^{n\Phi_{\mathrm{out}}^{\mathrm{sys}}(\eta)}.
\]
Theorem~\ref{prop:dense-inner-linear-exponent} gives
\[
p_h(\delta)
\le
n^{O(1)}\,2^{-nE_{\mathrm{in}}(\eta;\theta,\delta,\xi)}.
\]
Multiplying and using \eqref{eq:dense-dense-linear-gap},
\[
\EE[\enum_{\Cone,h}]\,p_h(\delta)
\le
n^{O(1)}\,2^{-\tau n}.
\]
There are only \(O(n)\) such weights \(h\), so summing over
\(h\in[\eta_0 n,(1-\theta)n]\) preserves the same exponential rate.
\end{proof}

\paragraph{Meaning of the criterion.}
Theorem~\ref{thm:dense-dense-linear-weight-criterion} isolates the real quantitative problem left by the baseline proof.
Once the low and intermediate weights are handled by the coarse argument, the only remaining analytic task is to show
that the outer exponent \(\Phi_{\mathrm{out}}^{\mathrm{sys}}\) sits below the dense-inner exponent
\(E_{\mathrm{in}}\) on the linear-weight window of interest.

\begin{theorem}[A concrete moderate-to-high-weight exponent gap]
\label{prop:dense-dense-highweight-gap}
Take
\[
\delta=0.12,
\qquad
\theta=0.005,
\qquad
\xi=8.
\]
Then direct evaluation of the explicit rate functions from
Corollary~\ref{cor:outer-dense-highweight-exponent} and
Theorem~\ref{prop:dense-inner-linear-exponent} shows that
\[
\Phi_{\mathrm{out}}^{\mathrm{sys}}(\eta)
\;
-\;
E_{\mathrm{in}}(\eta;0.005,0.12,8)
\le -0.13
\qquad
\text{for every }\eta\in[\eta_{\mathrm{crit}},0.99],
\]
where
\[
\eta_{\mathrm{crit}}:=1-\frac{1}{\sqrt2}.
\]
Consequently,
\[
\sum_{h=\lceil \eta_{\mathrm{crit}}n\rceil}^{\lfloor 0.99n\rfloor}
\EE[\enum_{\Cone,h}]\,p_h(0.12)
\le
n^{O(1)}\,2^{-0.13n}.
\]
\end{theorem}

\begin{proof}
The inner exponent is
\[
E_{\mathrm{in}}(\eta;0.005,0.12,8)
=
\min\!\left\{
\Psi_{\mathrm{run}}^{(w)}(\eta,0.005),\
10\cdot 0.12\left(1-H_2(1/10)\right)
\right\}.
\]
The explicit binomial term equals
\[
10\cdot 0.12\left(1-H_2(1/10)\right)\approx 0.6372.
\]
On \([\eta_{\mathrm{crit}},0.99]\), Corollary~\ref{cor:outer-dense-highweight-exponent} gives
\[
\Phi_{\mathrm{out}}^{\mathrm{sys}}(\eta)=H_2(\eta)-\frac12.
\]
A direct numerical evaluation of the explicit formulas on \([\eta_{\mathrm{crit}},0.99]\) then shows that the gap
\[
\Phi_{\mathrm{out}}^{\mathrm{sys}}(\eta)-E_{\mathrm{in}}(\eta;0.005,0.12,8)
\]
never exceeds \(-0.13\); the worst point occurs near \(\eta\approx 0.5\).
This verification is reproduced by the companion script
\texttt{scripts/verify\_dense\_claims.py}, which for the present parameters reports a worst sampled gap
of approximately \(-0.1372\) near \(\eta=0.5000\).
The final summation bound is then exactly Theorem~\ref{thm:dense-dense-linear-weight-criterion}.
\end{proof}

\begin{theorem}[A concrete twelve-percent dense+dense theorem]
\label{thm:dense-dense-twelve-percent}
Assume
\[
c>\frac{1}{\log_2(3/2)},
\qquad
\gamma>2.
\]
Then the systematic dense outer with memory
\[
M=\lceil c\log_2 n\rceil
\]
concatenated with the dense inner of memory
\[
m=\lceil \gamma\log_2 n\rceil
\]
satisfies
\[
\Pr[d_{\min}(\Csc)\le 0.12 n]=o(1).
\]
\end{theorem}

\begin{proof}
Set
\[
\delta=0.12,
\qquad
\kappa=\frac12,
\qquad
\rho=\frac14,
\qquad
\theta=0.005,
\qquad
\xi=8,
\qquad
h_0:=\lfloor \kappa\log_2 n\rfloor,
\]
and let \(\eta_{\mathrm{crit}}:=1-2^{-1/2}\).
Split the first-moment sum into four windows:
\[
S_{\mathrm{tiny}}
:=
\sum_{h=1}^{h_0}\EE[\enum_{\Cone,h}]\,p_h(0.12),
\]
\[
S_{\mathrm{low}}
:=
\sum_{h=h_0+1}^{\lfloor \eta_{\mathrm{crit}}n\rfloor}\EE[\enum_{\Cone,h}]\,p_h(0.12),
\]
\[
S_{\mathrm{lin}}
:=
\sum_{h=\lceil \eta_{\mathrm{crit}}n\rceil}^{\lfloor 0.99n\rfloor}\EE[\enum_{\Cone,h}]\,p_h(0.12),
\]
and
\[
S_{\mathrm{top}}
:=
\sum_{h=\lceil 0.99n\rceil}^{n}\EE[\enum_{\Cone,h}]\,p_h(0.12).
\]

For \(S_{\mathrm{tiny}}\), the proof of Corollary~\ref{cor:dense-dense-concrete-choice} already applies verbatim,
since \(0.12<1/8\), and gives \(S_{\mathrm{tiny}}=o(1)\).

For \(S_{\mathrm{low}}\), Corollary~\ref{cor:outer-dense-lowlinear-envelope} gives
\[
\EE[\enum_{\Cone,h}] \le n^{O(1)}(1+\sqrt2)^h
\qquad(h\le \eta_{\mathrm{crit}}n).
\]
Also, Corollary~\ref{cor:dense-inner-constant-geometric} applies with \(\rho=1/4\), \(\kappa=1/2\), and \(\gamma>2\),
so for all \(h\ge h_0+1\),
\[
p_h(0.12)\le 3\cdot 4^{-h}.
\]
Therefore
\[
S_{\mathrm{low}}
\le
n^{O(1)}\sum_{h=h_0+1}^{\lfloor \eta_{\mathrm{crit}}n\rfloor}
\left(\frac{1+\sqrt2}{4}\right)^h.
\]
Since \((1+\sqrt2)/4<1\), this is a geometric tail beginning at \(h_0=\tfrac12\log_2 n+O(1)\), hence
\[
S_{\mathrm{low}}
\le
n^{-\frac12\log_2(4/(1+\sqrt2))+o(1)}
=
o(1).
\]

For \(S_{\mathrm{lin}}\), Theorem~\ref{prop:dense-dense-highweight-gap} gives
\[
S_{\mathrm{lin}}\le n^{O(1)}2^{-0.13n}=o(1).
\]

For \(S_{\mathrm{top}}\), use only the outer exponent and the trivial bound \(p_h(0.12)\le 1\).
If \(h=\lfloor \eta n\rfloor\) with \(\eta\in[0.99,1]\), then Corollary~\ref{cor:outer-dense-highweight-exponent} gives
\[
\Phi_{\mathrm{out}}^{\mathrm{sys}}(\eta)=H_2(\eta)-\frac12
\le
H_2(0.99)-\frac12
< -0.41.
\]
Hence Theorem~\ref{thm:outer-dense-linear-exponent} gives
\[
\EE[\enum_{\Cone,h}]
\le
n^{O(1)}2^{-0.41n},
\]
uniformly over \(h\in[0.99n,n]\), and so
\[
S_{\mathrm{top}}
\le
n^{O(1)}2^{-0.41n}
=
o(1).
\]

Combining the four windows proves
\[
\Pr[d_{\min}(\Csc)\le 0.12 n]
\le
S_{\mathrm{tiny}}+S_{\mathrm{low}}+S_{\mathrm{lin}}+S_{\mathrm{top}}
=
o(1).
\]
\end{proof}

\subsection{What must be sharpened next}
\label{sec:integration-dense-dense-block}

The dense proof is no longer merely existential.
The baseline theorem proves linear distance in full generality, and
Theorem~\ref{thm:dense-dense-twelve-percent} already certifies one explicit distance constant.
What remains is to push that certified constant upward toward the finite-length behavior seen in the dense enumerator.
At the analytic level, there are still two concrete bottlenecks.

\paragraph{Outer bottleneck.}
Corollary~\ref{lem:outer-dense-spectrum} gives only a global geometric envelope
\[
\EE[\enum_{\Cone,h}] \le \beta_{\mathrm{out}}^h.
\]
This is enough for the baseline theorem, but it is far too coarse for near-GV claims.
The new explicit outer exponent already replaces it in the concrete twelve-percent theorem.
To move further, one needs a sharper outer spectrum exponent in the moderate-weight regime,
not merely a better global base \(\beta_{\mathrm{out}}\).

\paragraph{Inner bottleneck.}
On the inner side, Corollary~\ref{cor:dense-inner-constant-geometric} replaces the detailed run-aware behavior by a
single geometric factor \(\rho^h\).
Again, this is perfect for the low-linear splice used in Theorem~\ref{thm:dense-dense-twelve-percent},
but it loses the linear-weight exponent information needed for stronger distance constants.
The natural next sharpening is to use the full dense-inner exponent more aggressively below the current
\(\eta_{\mathrm{crit}}\) splice point.

\paragraph{Current sharp obstruction.}
The present paper already certifies \(d_{\min}(\Csc)\ge 0.12n\) with high probability.
The remaining blocker is therefore no longer linear distance itself, but rather the quantitative gap to the
enumerator-guided curve.
At the moment, the proof changes gears at \(\eta_{\mathrm{crit}}=1-2^{-1/2}\):
below that point it uses the low-linear outer envelope together with the inner geometric tail, while above that point
it uses the explicit exponent comparison.
The next theorem should remove or lower that splice by comparing the explicit outer and inner exponents on a wider
moderate-weight window.

\paragraph{Role of the enumerator.}
The exact enumerator is therefore not the proof itself; it is the calibration tool that tells us where sharper
analytic work is worth spending effort.
In particular, it suggests which \((\sigma,k)\) regimes track the random-linear benchmark, and hence which parameter
windows should be targeted when replacing the coarse pair \((\beta_{\mathrm{out}},\rho)\) by sharper exponents.

\subsection{How the enumerator enters}
\label{sec:integration-dense-dense-enumerator}

The proof above is deliberately coarse.
For concrete parameter choices, the exact concatenated enumerator is a calibration tool:
it suggests which \((\sigma,k)\) regimes are promising, which low-weight windows actually dominate,
and where the analytic proof is losing the most.
In the main body we only use analytic bounds, but the exact enumerator can be recorded in an appendix and used to guide
the sharpening of \eqref{eq:dense-dense-parameter-conditions}.

\subsection{Dense+dense estimates (working save point)}
\label{sec:integration-dense-dense-estimates}

This subsection is a \emph{working estimates section}, not a theorem section.
Its purpose is to record the current model we use to think about the dense+dense regime before the corresponding
finite-length statements are turned into fully rigorous propositions.
The guiding principle is still the first moment:
\begin{equation}
\label{eq:dense-dense-working-mu}
\mu_{n,\sigma}(\delta)
\;:=\;
\EE\!\left[\#\{c\neq 0:\ \wt(c)\le \delta n\}\right]
\;\approx\;
\sum_{h=1}^{n}
A_h^{\mathrm{out}}\,
p_h^{\mathrm{in}}(\delta),
\end{equation}
where \(A_h^{\mathrm{out}}\) is the expected number of outer codewords of intermediate weight \(h\), and
\(p_h^{\mathrm{in}}(\delta)\) is the inner low-tail probability from a uniformly random weight-\(h\) slice input.
Whenever \(\mu_{n,\sigma}(\delta)<1\), Markov's inequality gives the useful surrogate
\[
\Pr[d_{\min}(\Csc)\le \delta n]\le \mu_{n,\sigma}(\delta).
\]

\paragraph{Why we reparameterize by \(\sigma=\lceil \log_2 k\rceil+c\).}
For the finite-length construction suggested by the enumerator, the natural design parameter is not raw \(\sigma\), but
the offset
\[
\sigma=\lceil \log_2 k\rceil+c,
\qquad c\in\{0,1,2,\dots\}.
\]
This is motivated by the small-weight formulas themselves.
For fixed small \(h\), the exact outer low-weight contribution behaves like a polynomial in \(\sigma\) times
\(k2^{-\sigma}\), hence mainly as a function of the offset \(c\).
Likewise, in the inner small-weight bound the OFF term appears through \(L2^{-m}\), and when \(m\) is identified with
\(\sigma\) and \(L\asymp n\asymp k\), this again behaves roughly like a constant times \(2^{-c}\).
So the first-order question is:
\emph{for a target \(\delta\), which offset \(c\) is enough?}

\paragraph{The current finite-\(n\) model.}
The working evaluator still splits \eqref{eq:dense-dense-working-mu} into three windows, but the philosophy has
become more focused than when this section started.

\medskip
\noindent
\textbf{1. Linear weights are no longer the mystery.}
For \(h=\eta n\) with \(\eta\ge \eta_{\mathrm{crit}}\), we use the explicit exponent comparison already proved in the
main text:
\[
A_h^{\mathrm{out}}\,p_h^{\mathrm{in}}(\delta)
\;\lesssim\;
n^{O(1)}\,2^{n(\Phi_{\mathrm{out}}^{\mathrm{sys}}(\eta)-E_{\mathrm{in}}(\eta;\theta,\delta,\xi))}.
\]
This window is currently not the dominant obstruction in the parameter ranges suggested by the enumerator.
The dense+dense proof problem is now much more localized than that.

\medskip
\noindent
\textbf{2. The outer side is largely understood at low weight.}
For the tiny-weight window we use the actual systematic-banded small-weight law suggested by the enumerator rather than
the earlier convolutional proxy:
\begin{equation}
\label{eq:dense-dense-working-outer-small}
A_h^{\mathrm{out,exact}}
\;=\;
k\cdot \frac{\binom{\sigma+1}{h-1}}{2^{\sigma+1}}
\;+\;
\sum_{\ell=2}^{k}(k-\ell+1)\frac{\binom{2\ell+\sigma-2}{h-2}}{2^{\ell+\sigma}}.
\end{equation}
Once this exact low-weight outer law is used, the dominant uncertainty shifts back onto the inner tiny/intermediate
regime rather than the outer count itself.

\medskip
\noindent
\textbf{3. The inner side should be built around exact late placement, not raw run tails.}
The old pointwise run-tail picture
\[
\min_{1\le r\le h}
\left(
\mathrm{RunTail}_n(h,r)+(L2^{-\sigma})^r+r\,b_{\delta,\xi}(n)
\right)
\]
is still useful as a baseline, but it is no longer the conceptual center.
The current philosophy is:
\begin{itemize}
\item the one-run base factor \(q_1(\delta)\approx 2\delta\) is now analytically anchored by
Corollary~\ref{cor:dense-inner-single-run};
\item the conditioned late-start power \(q_1(\delta)^s\) is analytically anchored by
Corollary~\ref{cor:first-run-late-qbase};
\item the exact low-weight core should therefore be ``exact late placement plus a correction,'' not a direct promotion
of the old \((L2^{-m})^r\) law.
\end{itemize}

\paragraph{General philosophy of the current dense model.}
The current working model is not trying to guess the minimum distance curve from scratch. It is trying to expose the
\emph{smallest theorem-shaped object that still matches the exact overlap data}.
At this point the philosophy is:
\begin{enumerate}
\item use exact or proved combinatorics whenever possible;
\item isolate the first place where the proof stops matching the exact tables;
\item treat the residual there as the real theorem target, rather than patching the whole first moment globally.
\end{enumerate}
This is why the current dense save point is organized around isolated low-weight slices and exact early-start
statistics, not around increasingly elaborate global fit parameters.

\paragraph{Current evidence lanes.}
The finite-\(n\) evidence is now best read as a bracket rather than a single claim.
At \(k=2^{20}\), \(\delta=0.12\), and \(n=2k+\sigma\), the current optimistic / central / pessimistic lanes cross
below \(1\) near
\[
c_{\mathrm{opt}}\approx 5,\qquad
c_{\mathrm{ctr}}\approx 6,\qquad
c_{\mathrm{pess}}\approx 8.
\]
On the sweep \(k=2^{10},2^{12},\dots,2^{20}\) at \(\delta=0.12\), the first offsets with estimated first moment below
\(1\) are currently
\[
\begin{array}{c|cccccc}
\log_2 k & 10 & 12 & 14 & 16 & 18 & 20\\
\hline
c_{\mathrm{opt}}  & 5 & 4 & 4 & 5 & 5 & 5\\
c_{\mathrm{ctr}}  & 4 & 4 & 4 & 5 & 6 & 6\\
c_{\mathrm{pess}} & 4 & 5 & 6 & 7 & 7 & 8
\end{array}
\]
and the current theorem-shaped safe lane still sits near \(c\approx 6\text{--}7\) after a modest explicit proof tax.
These numbers are not theorem statements. Their role is only to tell us which proof targets are worth pursuing.

\paragraph{What is solid versus what is still heuristic.}
By now several parts of the picture are genuinely solid:
\begin{itemize}
\item the single-run law \(p_1^{\mathrm{in}}(\delta)\approx 2\delta\);
\item the conditioned late-placement base
\[
\sum_{s=1}^{h}
\pi_{N,h}(s)\,
\frac{\binom{L-h+1}{s}}{\binom{N-h+1}{s}};
\]
\item the fact that fixed tiny weights are overwhelmingly supported on the isolated slice \(R(U)=h\);
\item the fact that deterministic short-gap merging is not the dominant large-\(k\) correction in the
\(h\approx 5\text{--}10\) window.
\end{itemize}
What remains heuristic is the \emph{correction around that proved base}. The key change since earlier drafts is that
this correction is no longer an amorphous global fit. It is now concentrated on the isolated slice and on explicit
early-start statistics there.

\paragraph{Current proof-oriented picture.}
The isolated slice \(R(U)=h\) is now the natural theorem core.
On that slice we have, in increasing order of sharpness:
\begin{itemize}
\item the exact late-placement base
\[
\PP[R(U)=h]\cdot \frac{\binom{L-h+1}{h}}{\binom{N-h+1}{h}};
\]
\item exact order-statistic control from Corollary~\ref{cor:isolated-order-late};
\item the exact early-count law from Corollary~\ref{cor:isolated-early-count};
\item the exact expected early-pair count from Corollary~\ref{cor:isolated-early-pair-mean}.
\end{itemize}
The important negative results are now just as useful as the positive ones:
\begin{itemize}
\item the one-episode correction is much too weak by itself;
\item a uniform per-early-start penalty \(\tau\) already tracks the isolated slice fairly well locally, but is too loose
as a large-\(k\) theorem lane;
\item a position-sensitive one-run penalty
\[
\tau_y(\delta)\approx (N-y+1)2^{-m}+\PP[\Binomial(N-y+1,\tfrac12)\le \delta N]
\]
works well through roughly \(h\le 4\), but then starts to undershoot in the battleground window \(h\approx 5\text{--}7\).
\end{itemize}
So the remaining problem is no longer ``find a better global heuristic.'' It is to explain the residual after the
position-sensitive one-run base.

\paragraph{The current best residual target.}
At the moment the cleanest residual statistic is the hypergeometric early-pair count on the isolated slice.
After factoring out the position-sensitive one-run base, the exact overlap tables are well described by an affine law
in
\[
\EE\!\left[\binom{J_L(U)}{2}\,\middle|\, \wt(U)=h,\ R(U)=h\right]
\;=\;
\binom{h}{2}\frac{\binom{N-L}{2}}{\binom{N-h+1}{2}}.
\]
More specifically, on the current overlap set the residual coefficients themselves vary smoothly with
\[
q_1=\frac{2\lfloor \delta N\rfloor}{N}
\qquad\text{and}\qquad \frac{1}{m},
\]
which makes this look much closer to a real analytic theorem target than to a one-off regression artifact.
This should still be read as a working estimate, not as a proved law, but it is now the most principled target we have.

\paragraph{Current save point.}
The dense+dense picture is now the following.
\begin{itemize}
\item The outer side is largely under control in the tiny-weight regime.
\item The linear-weight window is no longer the main bottleneck.
\item The proof should be organized around exact late placement on the isolated slice.
\item The next theorem ingredient is an explicit multi-episode correction built on top of the position-sensitive
one-run base, most naturally through the early-start hypergeometric law and especially its early-pair statistic.
\end{itemize}
This should be read only as the current state of the model, not as a proved theorem. Its role is to preserve the
current proof philosophy and identify the narrow next step.

\paragraph{The theorem-shaped inner target suggested by the model.}
The current computational evidence suggests that the right next proof target is not an arbitrary refinement of
\eqref{eq:dense-inner-run-opt-envelope}, but an \emph{exact late-placement base plus a controlled correction}. Writing
\[
\pi_{N,h}(s)
\;:=\;
\PP[\#\mathrm{runs}(U)=s]
\;=\;
\frac{\binom{h-1}{s-1}\binom{N-h+1}{s}}{\binom{N}{h}},
\]
for \(U\) uniform on the weight-\(h\) slice, the natural target inequality is
\begin{equation}
\label{eq:dense-dense-working-pairwise-target}
p_h^{\mathrm{in}}(\delta)
\;\le\;
\sum_{s=1}^{h}
\pi_{N,h}(s)\,
\frac{\binom{L-h+1}{s}}{\binom{N-h+1}{s}}
\;+\; \mathrm{Corr}_{N,h}(\delta),
\end{equation}
where \(L=\lceil 2\delta N\rceil\) and the correction term \(\mathrm{Corr}_{N,h}(\delta)\) should be negligible on the
tiny/intermediate window.
The point is not that \eqref{eq:dense-dense-working-pairwise-target} is proved already, but that it is now a sharply
formulated analytic target: the late-placement base itself is already proved, and what remains is to upper-bound the
correction \(\mathrm{Corr}_{N,h}(\delta)\) coming from low-weight outcomes that are \emph{not} explained by pure
late placement. The present overlap checks suggest that this correction is a genuine low-order term in the dominant
tiny/intermediate window, rather than evidence for a qualitatively different small-weight law. The older pairwise
run-mixture lanes should now be read only as structured candidate models for \(\mathrm{Corr}_{N,h}(\delta)\), not as
the theorem core itself.

\paragraph{Why the isolated slice is the right core.}
At the smallest fixed weights, Corollary~\ref{cor:run-count-full} suggests that one should separate
\(\mathrm{Corr}_{N,h}(\delta)\) into:
\begin{itemize}
\item an isolated-run correction on the dominant \(R(U)=h\) slice, and
\item a lower-order adjacency correction from the event \(R(U)<h\).
\end{itemize}
That looks like the cleanest route if we later try to turn the current finite-\(h\) evidence into an actual proof.
Numerically, the exact isolated-slice base
\[
\PP[R(U)=h]\cdot \frac{\binom{L-h+1}{h}}{\binom{N-h+1}{h}}
\]
already matches the exact one-run slice exactly and stays within a few bits through the dominant tiny weights on the
current overlap tables. The remaining residual on that isolated slice is smooth and monotone enough that it looks
like a real correction term, not a failure of the late-placement core itself.

\paragraph{A concrete isolated-slice working target.}
The simplest working version of that correction is
\begin{equation}
\label{eq:dense-dense-working-isolated-target}
p_h^{\mathrm{in}}(\delta)
\;\le\;
\PP[R(U)=h]\,
\frac{\binom{L-h+1}{h}}{\binom{N-h+1}{h}}\,
2^{\beta_\delta \binom{h}{2}}
\;+\;
\PP[R(U)<h]
\;+\;
\mathrm{Err}^{\mathrm{iso}}_{N,h}(\delta),
\end{equation}
with \(\beta_\delta\) small on the tiny-weight window and \(\mathrm{Err}^{\mathrm{iso}}_{N,h}(\delta)\) ideally negligible
there. This isolates the two real sources of loss that still have to be controlled: residual interaction between
otherwise isolated early episodes, and the lower-order adjacency event \(R(U)<h\).

\paragraph{What the exact isolated statistics buy us.}
Corollary~\ref{cor:isolated-order-late} gives a rigorous next rung above the late-placement base: on the isolated
slice one can count exactly the event that the \(r\)-th isolated start lies inside the final \(L\) positions. This is
the clean theorem object behind statements like ``there is at most one genuinely early isolated impulse.'' Current
numerics show that the \(r=2\) version is still too weak by itself to certify the observed small-weight regime, but it
is now a precise proved object rather than a heuristic guess.

Corollary~\ref{cor:isolated-early-count} sharpens that viewpoint further: on the isolated slice one can count exactly
how many isolated starts fall before the final \(L\) positions. Empirically, a model of the form
\[
\sum_{j=0}^h
\PP[J_L(U)=j\mid \wt(U)=h,\ R(U)=h]\tau^j
\]
with a single \(\tau\) already tracks the exact isolated slices within about a bit on the current overlap tables.
That is encouraging because it isolates the remaining analytic burden into a per-early-episode penalty.
At the same time, a uniform safe \(\tau\) is still too loose when inserted into the large-\(k\) first moment, so the
next useful theorem likely needs a position-sensitive or \(j\)-sensitive refinement rather than a single global
constant.

The natural position-sensitive version is to assign each early isolated start position \(y\le N-L\) its own
penalty \(\tau_y(\delta)\), and then average the product of those penalties over the exact isolated-start
configuration. On the isolated slice this would take the form
\[
\frac{1}{\binom{N-h+1}{h}}
\sum_{j=0}^h
\binom{L-h+1}{h-j}\,
e_j\!\bigl(\tau_1(\delta),\dots,\tau_{N-L}(\delta)\bigr),
\]
where \(e_j\) is the \(j\)-th elementary symmetric polynomial. The uniform-\(\tau\) model is exactly the special case
\(\tau_y\equiv \tau\), so the current overlap evidence can be read as saying that the right theorem probably lives one
level above the constant-\(\tau\) simplification, not in a completely different combinatorial regime.

The most natural first choice is to take \(\tau_y(\delta)\) from the one-run bound with remaining suffix length
\(N-y+1\), namely
\[
\tau_y(\delta)
\approx
(N-y+1)2^{-m}
\;+\;
\Pr[\Binomial(N-y+1,\tfrac12)\le \delta N].
\]
That position-sensitive one-run model does explain the isolated slice better than a global constant at the very
smallest weights, but it still undershoots once \(h\) reaches the \(5\text{--}7\) window driving the current
large-\(k\) transition. So the next theorem is still not just ``plug in the one-run law per early start''; it needs
an additional multi-episode interaction term on top of that exact early-count decomposition.

\paragraph{The current best residual law.}
One especially clean candidate is suggested by Corollary~\ref{cor:isolated-early-pair-mean}: after factoring out the
position-sensitive one-run model, the remaining overlap gap is well tracked by an affine function of the exact
expected number of early-start pairs
\[
\EE\!\left[\binom{J_L(U)}{2}\,\middle|\, \wt(U)=h,\ R(U)=h\right]
=
\binom{h}{2}\frac{\binom{N-L}{2}}{\binom{N-h+1}{2}}.
\]
On the current overlap tables the residual coefficients are well fit by
\[
a(q_1,m)\approx -1.1434 + 6.1146\,q_1 - \frac{15.4212}{m},
\]
\[
b(q_1,m)\approx 0.4950 - 1.0649\,q_1 - \frac{0.7693}{m},
\]
so the residual takes the working form
\[
\text{gap bits}
\approx
a(q_1,m)
+
b(q_1,m)\,
\EE\!\left[\binom{J_L(U)}{2}\,\middle|\, \wt(U)=h,\ R(U)=h\right].
\]
This does not yet prove the right theorem, but it is a much sharper target than a free global fit: it says the next
correction likely lives on the early-start hypergeometric law itself, not on the full run count.

The new one-episode lemma confirms that \(\mathrm{Corr}_{N,h}(\delta)\) cannot be controlled by looking only at the
first early ON-episode: that route is structurally valid but numerically much too loose. Hence the real target is a
genuine multi-episode correction, probably pairwise at minimum.

\paragraph{What must be improved next.}
The outer small-weight formula \eqref{eq:dense-dense-working-outer-small} has now largely been replaced, for
finite-\(n\) estimation purposes, by the exact low-weight systematic banded law suggested by the enumerator.
The remaining gap is therefore mostly on the inner side:
even after sharpening the outer tiny-weight treatment, overlap checks against the exact enumerator still indicate that
the present small-weight models remain somewhat optimistic in the small/intermediate-weight regime.
So the next refinement target is not the bulk linear window, but rather a sharper and still explicit description of the
inner small-weight law for roughly \(h\approx 5\text{--}12\), ideally upgrading the pairwise run-mixture picture from a
working estimate to a finite-\(n\) bound. The current best proof-oriented subproblem is therefore the isolated slice
\(R(U)=h\): prove a controlled correction around the exact late-placement base there, and then add the lower-order
adjacency remainder coming from \(R(U)<h\). Once that is in place, the next task is a cleaner interpolation from that
tiny/intermediate regime to the
linear-weight exponent window.
